\documentclass[11pt]{article}

\usepackage[utf8]{inputenc}
\usepackage[T1]{fontenc}
\usepackage{lmodern}
\usepackage[margin=1in]{geometry}
\usepackage{amsmath,amssymb}
\usepackage{microtype}
\usepackage{enumitem}
\usepackage{booktabs}
\usepackage{array}
\usepackage{xcolor}
\usepackage[colorlinks=true,linkcolor=black,citecolor=black,urlcolor=blue]{hyperref}
\usepackage[backend=biber,style=numeric,sorting=nyt]{biblatex}
\addbibresource{tdtsp_references.bib}

% Accent colors consistent with the companion slide deck
\definecolor{emphasisblue}{HTML}{9B2D63}
\definecolor{emphasisteal}{HTML}{2C7E8C}
\newcommand{\emphasize}[1]{\textcolor{emphasisblue}{\textit{#1}}}
\newcommand{\emphteal}[1]{\textcolor{emphasisteal}{\textit{#1}}}

\newcommand{\bench}[1]{\textbf{#1}}

\title{\textbf{Time Dependency in the TDTSP benchmark:\\
Three Hybrid Quantum--Classical Approaches}}
\author{Alma Arevalo Loyola\\ \small Optimization Team, SuperQ}
\date{June 2026}

\begin{document}

\maketitle

\begin{abstract}
\noindent
The Time-Dependent Traveling Salesperson Problem (TD-TSP) asks for a minimum-cost
tour when the cost of traversing an edge depends on the time at which it is traversed.
The study "Benchmarking quantum-vs-classical approaches for Time-Dependent TSP,
obtains its tours by solving a \emph{static} TSP on a traffic-scaled cost matrix and
only afterwards evaluating the route under time-varying traffic. This simplification
sidesteps the defining difficulty of the problem: the cost of each edge should be
computed from the \emph{actual arrival time} at its origin node. This report proposes and
analyzes three hybrid quantum--classical approaches that reintroduce true time dependency
\emph{during} optimization. For each approach we give a full description, explain
precisely how time dependency is handled, and discuss its benefits, challenges, and
limitations. We close with a comparative summary.
\end{abstract}

\tableofcontents

\section{Background and motivation}
\label{sec:background}

\subsection{The TD-TSP}
In the classical TSP one seeks a minimum-cost Hamiltonian tour over $n$ nodes given a
fixed cost matrix. In the \emphasize{Time-Dependent TSP (TD-TSP)} each arc $(i,j)$ has a
cost that is a function of the departure time,
\[
  D_{ij}(t) \;=\; \text{travel time from $i$ to $j$ when departing $i$ at time $t$},
\]
so the cost of a tour depends on the order \emph{and} the timing of its visits: the
arrival time at each node determines the departure time used to evaluate the next arc.
Time-dependent routing is the realistic model for urban logistics, traffic-aware routing,
and dynamic supply chains \cite{gendreau2015time}. Both the TSP and the TD-TSP are
NP-hard, and the cascading dependence of arc costs on arrival times makes the TD-TSP
substantially harder to model and to solve than its static counterpart.

\subsection{The FIFO property}
A travel-time model is realistic only if it satisfies the
\emphteal{First-In-First-Out (FIFO)} property: leaving a node later can never lead to an
earlier arrival. The standard construction that guarantees FIFO represents each arc by a
stepwise speed function, from which a continuous piecewise-linear travel-time function
$D_{ij}(t)$ is derived; FIFO holds whenever the slope of $D_{ij}(t)$ exceeds $-1$
\cite{ichoua2003vehicle}. The parameters of any FIFO piecewise-linear model can be
recovered by solving a small linear system per arc \cite{ghiani2014note}. This FIFO
$D_{ij}(t)$ is the common foundation for all three approaches below.

\subsection{The gap to be closed}
The current paper replaces the function $D_{ij}(t)$ by a single
\emphasize{traffic multiplier} per departure window,
\[
  D_{ij}(t) \;=\; m_t\, D_{ij},
\]
where $D_{ij}$ is a baseline travel time and $m_t$ is a morning/evening/night multiplier
(e.g.\ $m_t = 1.66$ in the morning, $1.095$ at night). Because every arc is scaled by the
\emph{same} factor within a window, time only rescales the whole map uniformly: it never
changes which tour is best, so the optimizer effectively solves a static TSP. The three
approaches below differ precisely in \emph{how} they put the true, arrival-time-dependent
cost back into the optimization loop, while still exploiting a quantum solver where it is
most useful: on a small, well-isolated subproblem.

\section{Approach 1: Classical construction with quantum refinement}
\label{sec:approach1}

\subsection{Description}
This approach follows a \emphteal{construct-then-refine} pattern. First, the static
multiplier is replaced by genuine FIFO piecewise-linear functions $D_{ij}(t)$ for every
arc \cite{ichoua2003vehicle,ghiani2014note}. A classical \emphasize{time-dependent
insertion heuristic} then builds an initial tour: candidate insertions are evaluated by
propagating arrival times along the partial tour, so the construction already ``sees''
time-varying costs. Finally, the tour is improved locally: congested or expensive
\emphasize{critical segments} are extracted and each is re-optimized as a small
Quadratic Unconstrained Binary Optimization (QUBO) problem solved on a quantum annealer
or gate-based device. A segment of length $k$ is encoded as a $k$-city TSP-QUBO using the
standard Ising formulation \cite{lucas2014ising}, with arc costs fixed to the value of
$D_{ij}(t)$ in the relevant departure-time \emph{bucket}. The improved segment is spliced
back into the tour and the process iterates, much as classical hybrids use the annealer as
a within-route intensification operator \cite{feld2019hybrid,holliday2025advanced}.

\subsection{How it manages time dependency}
Time dependency enters in two places. (i) During \emph{construction}, arrival times are
propagated so that insertion costs use $D_{ij}(t)$ rather than a static matrix. (ii)
During \emph{refinement}, the departure window of each critical segment is discretized
into time buckets; within a bucket the arc costs are treated as constant, which yields a
well-posed QUBO whose coefficients nevertheless reflect the time-of-day cost of that
segment. Bucketing trades exactness for tractability: finer buckets approximate the true
$D_{ij}(t)$ more closely at the price of more frequent re-solves.

\subsection{Benefits}
\begin{itemize}[leftmargin=1.4em]
  \item \textbf{Modular and low-risk.} The classical construction guarantees a feasible,
        time-aware tour even if the quantum step contributes nothing; quantum is pure
        upside.
  \item \textbf{Small, hardware-friendly QUBOs.} Only short segments are sent to the
        device, keeping qubit counts and connectivity demands modest.
  \item \textbf{Mature components.} Time-dependent insertion heuristics and the TSP-QUBO
        encoding are well established \cite{gendreau2015time,lucas2014ising}.
\end{itemize}

\subsection{Challenges and limitations}
\begin{itemize}[leftmargin=1.4em]
  \item \textbf{Bucketing error.} Treating costs as constant within a bucket reintroduces
        a (smaller) version of the multiplier simplification; the route is only as
        time-accurate as the bucket resolution.
  \item \textbf{Local scope.} Refining isolated segments cannot fix a globally poor
        visiting order produced by the heuristic.
  \item \textbf{Re-evaluation cost.} After each splice, downstream arrival times shift, so
        the tour must be re-timed --- repeated many times this is non-trivial.
\end{itemize}

\section{Approach 2: CP master with quantum subproblem solvers}
\label{sec:approach2}

\subsection{Description}
Here a \emphasize{Constraint Programming (CP)} model is the global ``brain'' that owns the
schedule, and the quantum device is delegated only hard combinatorial sub-blocks. The CP
model represents visits as interval variables and enforces a \emphasize{time-dependent
no-overlap / travel constraint} so that the master always reasons about true arrival
times; such a constraint was introduced for urban delivery by
\citeauthor{aguiar2015time} \cite{aguiar2015time}. The customers are partitioned into
clusters (a cluster-first decomposition), each cluster-level tour is solved as a TSP-QUBO
on the quantum device \cite{feld2019hybrid,lucas2014ising}, and the CP master stitches the
cluster tours together while enforcing global timing and side constraints (time windows,
precedence, service times). This division of labor mirrors hybrid decomposition methods
in which clustering is classical and the per-cluster TSP is quantum
\cite{feld2019hybrid,holliday2025advanced}.

\subsection{How it manages time dependency}
Time dependency is owned entirely by the CP master: the time-dependent travel constraint
propagates arrival times exactly, and any candidate stitching that violates the FIFO
timing or a time window is pruned by constraint propagation. The quantum subproblems
themselves are static (a fixed small cost matrix per cluster), but because the master
decides \emph{when} each cluster is entered and \emph{how} the clusters are sequenced in
time, the overall solution respects true time dependency. This cleanly separates the
``timing logic'' (CP) from the ``combinatorial search'' (quantum).

\subsection{Benefits}
\begin{itemize}[leftmargin=1.4em]
  \item \textbf{Rich constraints for free.} CP excels exactly at scheduling-rich variants
        --- time windows, precedence, no-overlap --- that are awkward to encode directly in
        a QUBO \cite{aguiar2015time}.
  \item \textbf{Exact timing.} Arrival times are propagated exactly by the master, so the
        global solution is genuinely time-feasible.
  \item \textbf{Scalable decomposition.} Clustering keeps each quantum subproblem small
        while the master coordinates at scale \cite{feld2019hybrid}.
\end{itemize}

\subsection{Challenges and limitations}
\begin{itemize}[leftmargin=1.4em]
  \item \textbf{Decomposition quality.} Cluster boundaries are decided before timing is
        fully known; a poor clustering caps the achievable solution quality.
  \item \textbf{Master--subproblem interface.} Round-trips between the CP master and a
        cloud quantum device add latency and engineering complexity; minimizing the number
        of quantum calls is important \cite{holliday2025advanced}.
  \item \textbf{Static subproblems.} Within a cluster the cost is fixed, so intra-cluster
        time variation is approximated rather than modeled exactly.
\end{itemize}

\section{Approach 3: FMQA black-box optimization}
\label{sec:approach3}

\subsection{Description}
The third approach avoids expressing the time-dependent objective in closed QUBO form
altogether. The true time-dependent tour cost --- including cascading arrival times and any
realistic detail a simulator can capture --- is treated as an expensive
\emphasize{black-box} function $f(\mathbf{x})$ of a binary encoding $\mathbf{x}$ of the
route. The \emphasize{Factorization Machine with Quantum Annealing (FMQA)} method
\cite{kitai2020designing} then runs a surrogate-model loop:
\begin{enumerate}[leftmargin=1.6em]
  \item train a Factorization Machine (FM) surrogate
        $\hat f(\mathbf{x}) = w_0 + \sum_i w_i x_i + \sum_{i<j}\langle v_i,v_j\rangle x_i x_j$
        on the evaluated samples; being a quadratic polynomial, it maps directly to a QUBO;
  \item solve that QUBO on an annealer / Ising machine to propose a promising route;
  \item evaluate the proposed route with the true time-dependent cost;
  \item add the sample to the dataset, retrain, and repeat.
\end{enumerate}
For the TSP specifically, the encoding choice is decisive: Gray-style bit labeling that
maps similar routes to nearby bit strings markedly improves convergence and reduces the
qubit count relative to a naive one-hot encoding \cite{koshikawa2024efficient}.

\subsection{How it manages time dependency}
This approach handles time dependency in the most faithful way of the three: the black-box
evaluator can compute the exact arrival-time-dependent cost (or even call a full traffic
simulator), with no bucketing and no static sub-instances. Time dependency never has to be
linearized or embedded in the QUBO --- it lives entirely inside $f(\mathbf{x})$, while the
annealer only ever optimizes the cheap quadratic surrogate of that function.

\subsection{Benefits}
\begin{itemize}[leftmargin=1.4em]
  \item \textbf{Exact, arbitrary objective.} Any cost the simulator can compute --- exact
        FIFO costs, stochastic congestion, penalties --- is optimized as-is.
  \item \textbf{Sample efficiency.} The FM surrogate reduces the number of expensive
        black-box evaluations, which is the dominant cost \cite{kitai2020designing}.
  \item \textbf{Compact encodings.} Gray labeling lowers the bit count for TSP, easing
        hardware demands \cite{koshikawa2024efficient}.
\end{itemize}

\subsection{Challenges and limitations}
\begin{itemize}[leftmargin=1.4em]
  \item \textbf{Surrogate fidelity.} A quadratic FM may poorly approximate a highly
        non-convex tour-cost landscape, leading to slow or premature convergence.
  \item \textbf{Encoding sensitivity.} Performance depends strongly on the bit labeling;
        a bad encoding traps the search in local minima \cite{koshikawa2024efficient}.
  \item \textbf{Scalability.} Demonstrations remain at modest sizes; scaling to large $n$
        requires techniques such as sliding-window or domain-wall encodings and remains an
        open research question.
\end{itemize}

\section{Comparison and recommendations}
\label{sec:comparison}

Table~\ref{tab:compare} summarizes how the three approaches trade off fidelity of time
modeling against maturity and risk.

\begin{table}[h]
\centering
\renewcommand{\arraystretch}{1.35}
\begin{tabular}{>{\raggedright\arraybackslash}p{2.6cm}
                >{\raggedright\arraybackslash}p{3.5cm}
                >{\raggedright\arraybackslash}p{3.2cm}
                >{\raggedright\arraybackslash}p{3.2cm}}
\toprule
\textbf{Approach} & \textbf{Time dependency} & \textbf{Main benefit} & \textbf{Main limitation}\\
\midrule
1. Classical + quantum refinement
  & Approximate (time-bucketed segments)
  & Modular, low-risk, small QUBOs
  & Bucketing error; local scope \\
2. CP master + quantum subproblems
  & Exact at the master; static per cluster
  & Handles rich scheduling constraints
  & Clustering quality; interface latency \\
3. FMQA black-box
  & Exact (simulator-evaluated)
  & Most faithful objective
  & Surrogate fidelity; scalability \\
\bottomrule
\end{tabular}
\caption{Comparison of the three hybrid approaches.}
\label{tab:compare}
\end{table}

\paragraph{Recommendation.}
The three approaches are complementary rather than competing. Approach~1 is the natural
\emph{first step}: it is low-risk, reuses mature components, and immediately removes the
uniform-multiplier simplification at the segment level. Approach~2 is preferable when the
real instance carries \emph{rich side constraints} (time windows, precedence, service
times), where CP's exact timing is most valuable. Approach~3 is the most
\emph{research-forward}: it models time dependency exactly and is the right vehicle for
objectives that resist closed-form QUBO encoding, at the cost of greater uncertainty about
scaling. A pragmatic program would prototype Approach~1, adopt Approach~2 for
constraint-heavy deployments, and pursue Approach~3 as the longer-term path to genuinely
time-dependent quantum optimization. In all three, the guiding principle is the same and
matches current hardware reality: the quantum device solves only a small, well-isolated
subproblem, while a classical master guarantees that the route is optimized against real,
time-dependent costs.

\printbibliography[heading=bibintoc,title={References}]

\appendix

\section{QUBO formulation details}
\label{app:qubo}

\subsection{Static TSP as a QUBO}
Following the standard Ising formulation \cite{lucas2014ising}, encode an $N$-city tour
with binary variables $x_{v,p}\in\{0,1\}$, where $x_{v,p}=1$ iff city $v$ occupies
position $p$ in the tour ($v,p \in \{1,\dots,N\}$). The cost Hamiltonian is
\[
  H \;=\; \underbrace{A\sum_{v=1}^{N}\Big(1-\sum_{p=1}^{N}x_{v,p}\Big)^{2}}_{\text{each city once}}
       \;+\; \underbrace{A\sum_{p=1}^{N}\Big(1-\sum_{v=1}^{N}x_{v,p}\Big)^{2}}_{\text{each position filled}}
       \;+\; \underbrace{B\sum_{(u,v)} w_{uv}\sum_{p=1}^{N} x_{u,p}\,x_{v,p+1}}_{\text{tour length}},
\]
where $w_{uv}$ is the arc cost and indices on $p$ are taken modulo $N$. Feasibility of the
two penalty terms is guaranteed when the penalty weight dominates the largest possible
routing gain, e.g.\ $A > B\max_{u,v} w_{uv}$. The number of binary variables is $N^{2}$;
compact labelings can reduce this to $O(N\log N)$ \cite{koshikawa2024efficient}.

\subsection{Time-bucketed costs (Approach 1)}
For a critical segment with departure window $[t_0,t_1)$, partition the window into buckets
$\{B_1,\dots,B_m\}$ and, within bucket $B_b$, fix the arc cost to its representative value
$w_{uv}^{(b)} = D_{uv}(\tau_b)$ for a sample time $\tau_b \in B_b$. Solving the segment under
each bucket yields a piecewise-constant approximation of the true time-dependent segment
cost; finer buckets reduce the approximation error at the cost of more QUBO solves. The
segment objective reuses the routing term above with $w_{uv}\!\leftarrow\! w_{uv}^{(b)}$.

\subsection{FM surrogate as a QUBO (Approach 3)}
The Factorization Machine surrogate
\[
  \hat f(\mathbf{x}) \;=\; w_0 \;+\; \sum_{i} w_i x_i \;+\; \sum_{i<j} \langle v_i, v_j\rangle\, x_i x_j,
  \qquad \langle v_i,v_j\rangle = \sum_{f=1}^{k} v_{i,f}\,v_{j,f},
\]
is already quadratic in the binary vector $\mathbf{x}$, so it \emph{is} a QUBO: the linear
coefficients $w_i$ form the diagonal and the pairwise terms $\langle v_i,v_j\rangle$ the
off-diagonal entries. Minimizing $\hat f$ on an Ising machine proposes the next route to
evaluate with the true (Phase~0) time-dependent cost; the new sample updates the FM
parameters $(w_0, w_i, v_{i,f})$ and the loop repeats \cite{kitai2020designing}. The rank
$k$ controls surrogate expressiveness versus overfitting, and the binary encoding of routes
(e.g.\ Gray labeling) governs convergence \cite{koshikawa2024efficient}.

\end{document}
